\section{USAJMO 2019}
        \begin{enumerate}
        	\item (\textit{USAJMO 2019}) \textbf{Note:} For any geometry problem whose statement begins with an asterisk $(*)$, the first page of the solution must be a large, in-scale, clearly labeled diagram. Failure to meet this requirement will result in an automatic 1-point deduction.


 There are $a+b$ bowls arranged in a row, numbered $1$ through $a+b$, where $a$ and $b$ are given positive integers. Initially, each of the first $a$ bowls contains an apple, and each of the last $b$ bowls contains a pear.


 A legal move consists of moving an apple from bowl $i$ to bowl $i+1$ and a pear from bowl $j$ to bowl $j-1$, provided that the difference $i-j$ is even. We permit multiple fruits in the same bowl at the same time. The goal is to end up with the first $b$ bowls each containing a pear and the last $a$ bowls each containing an apple. Show that this is possible if and only if the product $ab$ is even.


 



\textbf{Solution}

\textbf{Claim:} If $a$ and $b$ are both odd, then the end goal of achieving $b$ pears followed by $a$ apples is impossible.

 \textbf{Proof:} Let $A_1$ and $P_1$ denote the number of apples and the number of pears in odd-numbered positions, respectively. We show that $A_1 - P_1$ is invariant. Notice that if $i$ and $j$ are odd, then $A_1$ and $P_1$ both decrease by 1, as one apple and one pear are both moved from odd-numbered positions to even-numbered positions. If $i$ and $j$ are even, then $A_1$ and $P_1$ both increase by 1.

 Because the starting configuration $aa\ldots ab\ldots b$ has $\left\lfloor \frac{a}{2} \right\rfloor + 1$ odd-numbered apples and $\left\lfloor \frac{b}{2} \right\rfloor$ odd-numbered pears, the initial value of $A_1 - P_1$ is $\left\lfloor \frac{a}{2} \right\rfloor - \left\lfloor \frac{b}{2} \right\rfloor + 1$. But the desired ending configuration has $\left\lfloor \frac{b}{2}\right\rfloor + 1$ odd-numbered pears and $\left\lfloor \frac{a}{2} \right\rfloor$ odd-numbered apples, so $A_1 - P_1 = \left\lfloor \frac{a}{2} \right\rfloor - \left\lfloor \frac{b}{2} \right\rfloor - 1$. As $A_1 - P_1$ is invariant, it is impossible to attain the desired ending configuration. $\blacksquare$

 \textbf{Claim:} If at least one of $a$ and $b$ is even, then the end goal of achieving $b$ pears followed by $a$ apples is possible.

 \textbf{Proof:} Without loss of generality, assume $a$ is even. If only $b$ is even, then we can number the bowls in reverse, and swap apples with pears. We use two inductive arguments to show that $(a,b) = (2,b)$ is achievable for all $b\ge 1$, then to show that $(a,b)$ is achievable for all even $a$ and all $b\ge 1$.

 \textit{Base case:} $(a,b) = (2,1)$. Then we can easily achieve the ending configuration in two moves, by moving the apple in bowl \#1 and the pear in bowl \#3 so that everything is in bowl \#2, then finishing the next move.

 \textit{Inductive step:} suppose that for $a \ge 2$ (even) apples and $b \ge 1$ pears, that with $a$ apples and $b$ pears, the ending configuration is achievable. We will show two things: i) achievable with $a$ apples and $b+1$ pears, and ii) achievable with $a+2$ apples and $b$ pears.

 i) Apply the process on the $a$ apples and first $b$ pairs to get a configuration $bb\ldots ba\ldots ab$. Now we will "swap" the leftmost apple with the last pear by repeatedly applying the move on just these two fruits (as $a$ is even, the difference $i-j$ is even). This gives a solution for $a$ apples and $b+1$ pears. In particular, this shows that $(a,b) = (2,b)$ for all $b\ge 1$ is achievable.

 ii) To show $(a+2, b)$ is achievable, given that $(a,b)$ is achievable, apply the process on the last $a$ apples and $b$ pears to get the configuration $aabbb\ldots baaa\ldots a$. Then, because we have shown that 2 apples and $b$ pears is achievable in i), we can now reverse the first two apples and $b$ pears.

 Thus for $ab$ even, the desired ending configuration is achievable. $\blacksquare$

 Combining the above two claims, we see that this is possible if and only if $ab$ is even. -scrabbler94

 



\textbf{Solution 2}

"If" part:Note that two opposite fruits can be switched if they have even distance. If one of $a$, $b$ is odd and the other is even, then switch $1$ with $a+b$, $2$ with $a+b-1$, $3$ with $a+b-2$... until all of one fruit is switched. If both are even, then if $a>b$, then switch $1$ with $a+1$, $2$ with $a+2$, $3$ with $a+3$... until all of one fruit is switched; if $a<b$ then switch $a$ with $a+b$, $a-1$ with $a+b-1$, $a-2$ with $a+b-2$... until all of one fruit is switched. Each of these processes achieve the end goal.

 "Only if" part:Assign each apple a value of $1$ and each pear a value of $-1$. At any given point in time, let $E$ be the sum of the values of the fruit in even-numbered bowls. Because $a$ and $b$ both are odd, at the beginning there are $\frac{a-1}{2}$ apples in even bowls and $\frac{b+1}{2}$ pears in even bowls, so at the beginning $E=\frac{a-b}{2}-1$. After the end goal is achieved, there are $\frac{a+1}{2}$ apples in even bowls and $\frac{b-1}{2}$ pears in even bowls, so after the end goal is achieved $E=\frac{a-b}{2}+1$. However, because two opposite fruits must have the same parity to move and will be the same parity after they move, we see that $E$ is invariant, which is a contradiction; hence, it is impossible for the end goal to be reached if $ab$ is odd.

 Stormersyle

 



\textbf{Solution 3}

We define a cycle to be when $2$ fruits in bowls $i$ and $j$ are swapped with eachother, provided that bowls $i$ contains an apple, bowl $j$ contains a pear, and the difference $i-j$ is even. 

 We claim that a cycle consists of only the valid moves described in the problem statement. Proof: We can move the apple and pear one step at a time to eventually reach the other's bowl(i.e. the apple moves to $i+1$ and the pear moves to bowl $j-1$ at the beginning) using only valid moves. By a simple parity argument(if $k$ moves have already been done, $i+k-(j-k)=i+j+2k=0+2k=2k=0 \textbf{mod 2}$), these moves are always valid, so we are done. $\blacksquare$

 We now split the "if" part of this proof into cases based on the parities of $a$ and $b$. 

 Case 1: $a$ even, $b$ oddIn this case we can simply cycle the fruit in bowl $k+1$ with the fruit in bowl $a+b-k$ for all integer $k$ in the range $[0,b-1]$ to achieve our desired goal. The difference $a+b-k-k-1=a+b-1$ is even, so the move is valid.

 Case 2: $b$ even, $a$ oddUsing the same cycle as the previous case works here.

 Case 3: $a$ even, $b$ evenHere, we can cycle the fruits in bowls $k$ and $a+b-k$ for all integer $k$ in the range $[1,b]$. Then, we finish by cycling the fruits in bowls $a$ and $a+b$. These are both valid moves by similar parity arguments as done in Case 1. 

 Since we have exhausted all possible cases for the "if" part of this proof, we have completed the if part. We will proceed with the "only if" part. 

 We let the number of apples in odd positions be $A_1$, and the number of pears in odd positions be $A_2$. Obviously $A_1-A_2$ is invariant when using valid moves only. We can now use proof by contradiction(i.e. showing that this invariant does not occur when $a, b \equiv 1 \text{ mod } 2$ and our desired goal is achieved) to find that we have that $0=2$, which is obviously false, so we have completed the proof for the "only if" part. 

 We are now done, since we have completed our proofs for the "if" and "only if" parts. $\blacksquare$

 



\textbf{Solution 4}

Begin by defining a $\textit{swap}$ as a set of $i - j$ legal moves taking an apple in $i \rightarrow j$ and the pear in $j \rightarrow i$. Now we begin by proving if $ab$ is even, then the final configuration is possible.

 WLOG let $b$ be even. Then there are two cases to consider. $a < b, a \ge b$ to provide valid constructions for. In the first case, consider the apple/pear pairs $(1,b+1), (2, b+2), \dots, (i, b+i), \dots, (a, b+a)$ and swap them. This results in the transformation, 
 \begin{align*} \underbrace{AA \cdots A}_{\text{a apples}}& \underbrace{PP \cdots P}_{\text{b pears}} \\ \underbrace{PP \cdots P}_{\text{a pears}}& \underbrace{PP \cdots P}_{\text{b-a pears}} \underbrace{AA \cdots A}_{\text{a apples}}  \\ \underbrace{PP \cdots P}_{\text{b pears}} &\underbrace{AA \cdots A}_{\text{a apples}}  \\ \end{align*}
as desired.

 Now consider the case where $a \ge b$. Also consider the pairing, $(a-b + 1, a + 1), (a - b + 2, a+2), \dots, (a - b + i, a+i), \dots, (a, a + b)$ and swap them. This results in the transformation, 
 \begin{align*} \underbrace{AA \cdots A}_{\text{a apples}} & \underbrace{PP \cdots P}_{\text{b pears}} \\ \underbrace{AA \cdots A}_{\text{a-b apples}}\underbrace{PP \cdots P}_{\text{b pears}}& \underbrace{AA \cdots A}_{\text{b apples}}. \end{align*}
Now just consider the leftmost $a$ bowls. There are $a-b$ apples, and then $b$ pears. We have effectively reduced $(a,b) \rightarrow (a-b, b)$. Reducing in this fashion recursively will lead to $(a \pmod{b}, b)$ which can be constructed with the $a<b$ case. This finishes the construction for $a \ge b$.

 Now assume $ab,a,b$ are odd. FTSOC, let this case satisfy the claim. Then, for each apple, define it's $\textit{position}$ as the sum of the distance from it to all other pairs. So for apple in bowl $i \le a$ originally, the position of $i$ is $\sum_{j = a+1}^{b} j - i .$

 Every time a legal move is made, the position of the apple decreases by $b - 1 + 2 = b+1$ (from the $b-1$ pears not part of the move and the decrease in $2$ by the corresponding pear). Also note that there needs to be exactly $ab$ moves in order to attain the final configurations. So the difference in sum of positions of the apples at the beginning and end is $ab(b+1)$. Mathematically,
 \begin{align*} \sum_{i = 1}^{a} \sum_{j = a+1}^{a+b} j - i  - \sum_{i = b + 1}^{b+a} \sum_{j = 1}^{b} j - i &= ab(b+1)\\ (\frac{ab(2a + b + 1)}{2} - \frac{ab(a + 1)}{2}) - (\frac{ab(b+1)}{2} - \frac{ab(a + 2b + 1)}{2}) &=  \\ ab(a+b) &= ab(b+1) \\ a &= 1. \end{align*}
A symmetric argument for the position of pears gives $b = 1$. Clearly the $a=b=1$ does not have a solution, therefore we have contradiction. $\square$

 


 


	\item (\textit{USAJMO 2019}) \textbf{Note:} For any geometry problem whose statement begins with an asterisk $(*)$, the first page of the solution must be a large, in-scale, clearly labeled diagram. Failure to meet this requirement will result in an automatic 1-point deduction.


 Let $\mathbb Z$ be the set of all integers. Find all pairs of integers $(a,b)$ for which there exist functions $f:\mathbb Z\rightarrow\mathbb Z$ and $g:\mathbb Z\rightarrow\mathbb Z$ satisfying 
 \[f(g(x))=x+a\quad\text{and}\quad g(f(x))=x+b\] for all integers $x$.


 



\textbf{Solution 1}

We claim that the answer is $|a|=|b|$.

 Proof:$f$ and $g$ are surjective because $x+a$ and $x+b$ can take on any integral value, and by evaluating the parentheses in different order, we find $f(g(f(x)))=f(x+b)=f(x)+a$ and $g(f(g(x)))=g(x+a)=g(x)+b$. We see that if $a=0$ then $g(x)=g(x)+b$ to $b=0$ as well, so similarly if $b=0$ then $a=0$, so now assume $a, b\ne 0$. 

 We see that if $x=|b|n$ then $f(x)\equiv f(0) \pmod{|a|}$, if $x=|b|n+1$ then $f(x)\equiv f(1)\pmod{|a|}$, if $x=|b|n+2$ then $f(x)\equiv f(2)\pmod{|a|}$... if $x=|b|(n+1)-1$ then $f(x)\equiv f(|b|-1)\pmod{|a|}$. This means that the $b$-element collection $\left\{f(0), f(1), f(2), ... ,f(|b|-1)\right\}$ contains all $|a|$ residues mod $|a|$ since $f$ is surjective, so $|b|\ge |a|$. Doing the same to $g$ yields that $|a|\ge |b|$, so this means that only $|a|=|b|$ can work.

 For $a=b$ let $f(x)=x$ and $g(x)=x+a$, and for $a=-b$ let $f(x)=-x$ and $g(x)=-x-a$, so  $|a|=|b|$ does work and are the only solutions, as desired.

 -Stormersyle

 



\textbf{Solution 2}

We claim that $f$ and $g$ exist if and only if $|a|=|b|$.

 \textbf{Only If:}

 For some fixed $j$, let $f(j)=k$.

 If $b=0$, then $g(k)=j$. Suppose $a\ne 0$. Then $f(j)=f(g(k))=k+a\ne k$, a contradiction. Thus, $a=0$. Similarly, if $a=0$, then $b=0$, satisfying $|a|=|b|$.

 Otherwise, $a,b\ne 0$. We know that $g(k)=g(f(j))=j+b$, $f(j+b)=f(g(k))=k+a$, $g(k+a)=j+2b$, and so on: $f(j+nb)=k+na$ and $g(k+na)=j+(n+1)b$ for $n\ge 0$.

 Consider the value of $g(k-a)$. Suppose $g(k-a)=j'\ne j$. Then $f(j')=k$ and $g(f(j'))=j+b\ne j'+b$, a contradiction. Thus, $g(k-a)=j$. We repeat with $f(j-b)$. Suppose $f(j-b)=k'-b\ne k-b$. Then $g(k'-b)=j$ and $f(g(k'-b))=k\ne k'$, a contradition. Thus, $f(j-b)=k-b$. Continuing, $g(k-2a)=j-a$, and so on: $f(j+nb)=k+na$ and $g(k+na)=j+(n+1)b$ now for all $n$.

 This defines $f(x)$ for all $x\equiv j\pmod{|b|}$ and $g(x)$ for all $x\equiv f(j)\pmod{|a|}$.

 This means that $x\equiv j\pmod{|b|}\implies f(x)\equiv f(j)\pmod{|a|}$, and $y\equiv f(j)\pmod{|a|}\implies g(y)\equiv j\pmod{|b|}$ which implies $f(x)\equiv f(j)\pmod{|a|}\implies x\equiv j\pmod{|b|}$.

 As a result, $f(x)$ maps each residue mod $|b|$ to a unique residue mod $|a|$, so $|a|\ge|b|$. Similarly, $g(x)$ maps each residue mod $|a|$ to a unique residue mod $|b|$, so $|b|\ge|a|$. Therefore, $|a|=|b|$.

 \textbf{If:}

 $|a|=|b|$ means that either $a=b$ or $a=-b$. $f(x)=x,g(x)=x+a$ works for the former and $f(x)=-x,g(x)=-x-a$ works for the latter, and we are done.

 


	\item (\textit{USAJMO 2019}) \textbf{Note:} For any geometry problem whose statement begins with an asterisk $(*)$, the first page of the solution must be a large, in-scale, clearly labeled diagram. Failure to meet this requirement will result in an automatic 1-point deduction.


 $(*)$  Let $ABCD$ be a cyclic quadrilateral satisfying $AD^2+BC^2=AB^2$. The diagonals of $ABCD$ intersect at $E$. Let $P$ be a point on side $\overline{AB}$ satisfying $\angle APD=\angle BPC$. Show that line $PE$ bisects $\overline{CD}$.


 



 



\textbf{Solution 1}

Let $PE \cap DC = M$. Also, let $N$ be the midpoint of $AB$.

 Note that only one point $P$ satisfies the given angle condition. With this in mind, construct $P'$ with the following properties:

 $AP' \cdot AB = AD^2 \quad \text{and} \quad BP' \cdot AB = CD^2$

 Claim: $P = P'$

 Proof:

 The conditions imply the similarities $ADP \sim ABD$ and $BCP \sim BAC$ whence $\measuredangle APD = \measuredangle BDA = \measuredangle BCA = \measuredangle CPB$ as desired. $\square$

 Claim: $PE$ is a symmedian in $AEB$

 Proof:

 We have 

 
 \[AP \cdot AB = AD^2 \iff AB^2 \cdot AP = AD^2 \cdot AB\]
 \[\iff \left( \frac{AB}{AD} \right)^2 = \frac{AB}{AP}\]
 \[\iff \left( \frac{AB}{AD} \right)^2 - 1 = \frac{AB}{AP} - 1\]
 \[\iff \frac{AB^2 - AD^2}{AD^2} = \frac{BP}{AP}\]
 \[\iff \left(\frac{BC}{AD} \right)^2 = \left(\frac{BE}{AE} \right)^2 = \frac{BP}{AP}\]

 as desired. $\blacksquare$

 Since $P$ is the isogonal conjugate of $N$, $\measuredangle PEA = \measuredangle MEC = \measuredangle BEN$. However $\measuredangle MEC = \measuredangle BEN$ implies that $M$ is the midpoint of $CD$ from similar triangles, so we are done. $\blacksquare$

 



\textbf{Solution 2}

By monotonicity, we can see that the point $P$ is unique. Therefore, if we find another point $P'$ with all the same properties as $P$, then $P = P'$

 Part 1) Let $N$ be a point on $\overline{AB}$ such that $AN\cdot AB = AD^2$, and $BN \cdot AB = BC^2$. Obviously $N$ exists because adding the two equations gives $AN\cdot AB + BN \cdot AB = AD^2 + BC^2 = AB^2$, which is the problem statement. Notice that converse PoP gives
 \[AN\cdot AB = AD^2 \implies \bigtriangleup ADN \sim \bigtriangleup ABD\]
 \[BN\cdot AB = AD^2 \implies \bigtriangleup CBN \sim \bigtriangleup ABC\]Therefore, $\angle AND = \angle ADB = \angle ACB = \angle CNB$, so $N$ does indeed satisfy all the conditions $P$ does, so $N = P$. Hence, $\bigtriangleup ADP \sim \bigtriangleup ABD$ and $\bigtriangleup CBP \sim\bigtriangleup ABC$.

 Part 2) Define $G$ as the midpoint of $CD$. Furthermore, create a point $X$ such that $DX || EC$ and $CX || ED$. Obviously $XCED$ must be a parallelogram. Now we set up for Jacobi's. The problem already gives us that $\angle APD = \angle CPB$, which is good for starters. Furthermore, $\bigtriangleup ADP \sim \bigtriangleup ABD$ tells us that
 \[\angle ADP = \angle ABD = \angle ACD = \angle XDC\]This gives us our second needed angle equivalence. Lastly, $\bigtriangleup CBP \sim\bigtriangleup ABC$ will give
 \[\angle BCP = \angle BAC = \angle BDC = \angle XCD\]which is our last necessary angle equivalence to apply Jacobi's. Finally, applying Jacobi's tells us that $AC$, $BD$, and $XP$ are concurrent $\implies$ $X$, $E$, $P$ collinear. Additionally, since parallelogram diagonals bisect each other, $X$, $G$, and $E$ are collinear, so finally we obtain that $P$, $E$, and $G$ are collinear, as desired.

 -jj\_ca888

 


	\item (\textit{USAJMO 2019}) \textbf{Note:} For any geometry problem whose statement begins with an asterisk $(*)$, the first page of the solution must be a large, in-scale, clearly labeled diagram. Failure to meet this requirement will result in an automatic 1-point deduction.


 $(*)$ Let $ABC$ be a triangle with $\angle ABC$ obtuse. The $A$\textit{-excircle} is a circle in the exterior of $\triangle ABC$ that is tangent to side $BC$ of the triangle and tangent to the extensions of the other two sides. Let $E, F$ be the feet of the altitudes from $B$ and $C$ to lines $AC$ and $AB$, respectively. Can line $EF$ be tangent to the $A$-excircle?


 



\textbf{Solution 1}

Instead of trying to find a synthetic way to describe $EF$ being tangent to the $A$-excircle (very hard), we instead consider the foot of the perpendicular from the $A$-excircle to $EF$, hoping to force something via the length of the perpendicular.  It would be nice if there were an easier way to describe $EF$, something more closely related to the $A$-excircle; as we are considering perpendicularity, if we could generate a line parallel to $EF$, that would be good.  

 So we recall that it is well known that triangle $AEF$ is similar to $ABC$.  This motivates reflecting $BC$ over the angle bisector at $A$ to obtain $B'C'$, which is parallel to $EF$ for obvious reasons.  

 Furthermore, as reflection preserves intersection, $B'C'$ is tangent to the reflection of the $A$-excircle over the $A$-angle bisector.  But it is well-known that the $A$-excenter lies on the $A$-angle bisector, so the $A$-excircle must be preserved under reflection over the $A$-excircle.  Thus $B'C'$ is tangent to the $A$-excircle.Yet for all lines parallel to $EF$, there are only two lines tangent to the $A$-excircle, and only one possibility for $EF$,  so $EF = B'C'$.  

 Thus as $ABB'$ is isoceles, 
 \[[ABC] = \frac{1}{2} \cdot AC \cdot BE = \frac{AC}{2} \cdot \sqrt{AB^2 - AE^2} = \frac{AC}{2} \cdot \sqrt{AB^2 - AB'^2} = \frac{AC}{2} \cdot \sqrt{AB^2 - AB^2} = 0,\] contradiction. -alifenix-

 



\textbf{Solution 2}

The answer is no.

 Suppose otherwise. Consider the reflection over the bisector of $\angle BAC$. This swaps rays $AB$ and $AC$; suppose $E$ and $F$ are sent to $E'$ and $F'$. Note that the $A$-excircle is fixed, so line $E'F'$ must also be tangent to the $A$-excircle.

 Since $BEFC$ is cyclic, we obtain $\measuredangle ECB = \measuredangle EFB = \measuredangle EF'E'$, so $\overline{E'F'} \parallel \overline{BC}$. However, as $\overline{EF}$ is a chord in the circle with diameter $\overline{BC}$, $EF \le BC$.

 If $EF < BC$ then $E'F' < BC$ too, so then $\overline{E'F'}$ lies inside $\triangle ABC$ and cannot be tangent to the excircle.

 The remaining case is when $EF = BC$. In this case, $\overline{EF}$ is also a diameter, so $BECF$ is a rectangle. In particular $\overline{BE} \parallel \overline{CF}$. However, by the existence of the orthocenter, the lines $BE$ and $CF$ must intersect, contradiction.

 



\textbf{Solution 3}

The answer is $\boxed{\text{no}}$.

 Suppose for the sake of contradiction that it is possible for $EF$ to be tangent to the $A$-excircle. Call the tangency point $T$, and let $S_1, S_2$ denote the contact points of $AB, AC$ with the $A$-excircle, respectively. Let $s$ denote the semiperimeter of $ABC$. By equal tangents, we have 
 \[ET = ES_2, FT = FS_1 \implies EF = ES_2+FS_2\]It is also well known that $AS_1 = AS_2 = \frac{s}{2}$, so 
 \[EF = ES_2+FS_2 = (AS_2-AE)+(AS_1-AF) = s-AE-AF \implies s=AE+AF+EF\]It is well known (by an easy angle chase) that $\triangle AEF \sim \triangle ABC$, so we must have the ratio of similitude is $2$. In particular, 
 \[AB=2 \cdot AE, AC=2 \cdot AF\]This results in 
 \[\angle ABE = 30^{\circ}, \angle CBF = 30^{\circ} \implies \angle EBC = 120^{\circ}\]which is absurd since $\triangle BEC$ is a right triangle. We reached a contradiction, so we are done. $\blacksquare$ ~ Mathscienceclass

 



\textbf{Solution 4}


                    \begin{figure}[H]
                        \centering
                        \includegraphics[width=0.5\textwidth]{Solution/USAJMO_I_2019_4_Solution_Figure_1.png}
                    \end{figure}
                    We claim that the answer is no. We proceed with contradiction. Suppose that $EF$ is indeed tangent to the a-excenter. Define the point of tangency to be $X$. Let $G$ to be the intersection of the a-excircle with the extension of $AB$ and $H$ to be the intersection of the a-excircle with the extension of $AC$. Define $I_a$ to be the a-excenter. It is a well-known fact that $I_a$ lies on the angle bisector of $\angle BAC$. For convenience, let $AB = c, AC = b, BC = a$, $\angle CAB = A, \angle CBA = B, \angle ACB = C$.Notice that by Power of a Point:
 \[FX = FG.\]
 \[XE = EH.\]Therefore, adding, we see that:
 \[FE = FX + XE = FG + EH.\]It would be rather nice if we could re-write $FG$. Indeed, we can:
 \[FG = AG - FG = \frac{r_a}{\tan \left( \frac{A}{2} \right)} - b \cos (A)\]and similarly for $FH$:
 \[FH = AH - AF = \frac{r_a}{\tan \left( \frac{A}{2} \right)} - c \cos (A)\]We now seek to re-write $FE$. Using the law of cosines:
 \[FE^2 = AF^2 + AE^ 2- 2AF \cdot AE \cos (A) = \left(c\cos (A) \right)^2 + \left(b\cos (A) \right)^2 -  2 (bc \cos ^2 (A)) \cos (A) = \cos ^2 (A) (c^2 + b^2 - 2bc \cos (A)) = \cos ^2 (A) \cdot a^2\]Therefore, 
 \[FE = a \cos (A)\]Putting this all together, we see that:
 \[a \cos (A) = \frac{r_a}{\tan \left(\frac{a}{2} \right)} - b \cos (A) + \frac{r_a}{\tan \left(\frac{a}{2} \right)} - c \cos (A) \Longleftrightarrow\]
 \[(a+b+c) \cos (A) = \frac{2r_a}{\tan \left(\frac{a}{2} \right)}\]Now, we seek to write $\tan \left(\frac{a}{2} \right)$ in terms of the side lengths of the triangle. Notice that:
 \[2\cos ^2 \left(\frac{A}{2} \right)-1 = \cos (A) \Longleftrightarrow\]
 \[2\cos ^2 \left(\frac{A}{2} \right) = 1 + \frac{b^2 + c^2-a^2}{2bc} = \frac{b^2 + 2bc + c^2 - a^2}{2bc} = \frac{(b+c)^2-a^2}{2bc} = \frac{(a+b+c)(b+c-a)}{2bc} = \frac{(2s)(2s-2a)}{2bc}\]where $s$ is the semi-perimeter. We get that:
 \[\cos \left(\frac{A}{2} \right) = \sqrt{\frac{s(s-a)}{bc}}\]Using the fact that $\cos^2 \left(\frac{A}{2} \right) + \sin^2 \left(\frac{A}{2} \right) = 1$, we have that:
 \[\sin^2 \left(\frac{A}{2} \right) = 1-\frac{b^2 + 2bc + c^2 - a^2}{4bc} = \frac{(b-c)^2 - a^2}{4bc} = \frac{(2s-2c)(2s-2b)}{4bc}\]Therefore, 
 \[\tan \left( \frac{A}{2} \right) = \sqrt{ \frac{(s-b)(s-c)}{s(s-a)} }\]Returning to our original problem:
 \[\frac{2r_a}{\sqrt{ \frac{(s-b)(s-c)}{s(s-a)} }} = (a+b+c) \cos (A)\]It is a well-known fact that $r_a = \sqrt{\frac{s(s-b)(s-c)}{s-a}}$, so:
 \[\frac{2\sqrt{\frac{s(s-b)(s-c)}{s-a}}}{\sqrt{ \frac{(s-b)(s-c)}{s(s-a)} }} = 2s \cos (A)\] which implies that:
 \[\cos (A) =1 \Longleftrightarrow A = 0\]which is a contradiction. Hence, our original assumption that $EF$ is tangent to the a-excircle is incorrect. $\blacksquare$~AopsUser101

 


	\item (\textit{USAJMO 2019}) \textbf{Note:} For any geometry problem whose statement begins with an asterisk $(*)$, the first page of the solution must be a large, in-scale, clearly labeled diagram. Failure to meet this requirement will result in an automatic 1-point deduction.


 Let $n$ be a nonnegative integer. Determine the number of ways that one can choose $(n+1)^2$ sets $S_{i,j} \subseteq \{1,2,...,2n\}$, for integers $i,j$ with $0 \le i,j \le n$ such that:


 $\bullet$ for all $0 \le i,j \le n$, the set $S_{i,j}$ has $i+j$ elements; and


 $\bullet$ $S_{i,j} \subseteq S_{k,l}$ whenever $0 \le i \le k \le n$ and $0 \le j \le l \le n$


 



\textbf{Solution 1}

Note that there are $(2n)!$ ways to choose $S_{1, 0}, S_{2, 0}... S_{n, 0}, S_{n, 1}, S_{n, 2}... S_{n, n}$, because there are $2n$ ways to choose which number $S_{1, 0}$ is, $2n-1$ ways to choose which number to append to make $S_{2, 0}$, $2n-2$ ways to choose which number to append to make $S_{3, 0}$... After that, note that $S_{n-1, 1}$ contains the $n-1$ in $S_{n-1. 0}$ and 1 other element chosen from the 2 elements in $S_{n, 1}$ not in $S_{n-1, 0}$ so there are 2 ways for $S_{n-1, 1}$. By the same logic there are 2 ways for $S_{n-1, 2}$ as well so $2^n$ total ways for all $S_{n-1, j}$, so doing the same thing $n-1$ more times yields a final answer of $(2n)!\cdot 2^{n^2}$.

 -Stormersyle

 



\textbf{Solution 2}

There are $\frac{(2n)!}{2^n}$ ways to choose $S_{0,0}, S_{1,1} .... S_{n,n}$. Since, there are $\binom{2n}{2}$ ways to choose $S_{1,1}$, and after that, to generate $S_{2,2}$, you take $S_{1,1}$ and add 2 new elements, getting you $\binom{2n-2}{2}$ ways to generate $S_{2,2}$. And you can keep going down the line, and you get that there are $\frac{(2n)!}{2^n}$ ways to pick $S_{0,0}, S_{1,1} ... S_{n,n}$ Then we can fill out the rest of the gird. First, let’s prove a lemma.

 \textbf{Lemma} \textbf{Claim:} If we know what $S_{a,b}$ is and what $S_{a+1,b+1}$ is, then there are 2 choices for both $S_{a,b+1}$ and $S_{a+1,b}$. 

 \textbf{Proof:} Note $a \leq a+1$ and $b \leq b+1$, so $S_{a,b}\subseteq S_{a+1,b+1}$. Let $A$ be a set that contains all the elements in $S_{a+1,b+1}$ that are not in $S_{a,b}$. $A = S_{a+1,b+1} - S_{a,b}$. We know $S_{a+1,b+1}$ contains total $a+b+2$ elements. And $S_{a,b}$ contains total $a+b$ elements. That means $A$ contains only 2 elements since $(a+b+2)-(a+b) = 2$. Let’s call these 2 elements $m, n$. $A = \{m, n\}$. $S_{a,b+1}$ contains 1 elements more than $S_{a,b}$ and 1 elements less than $S_{a+1,b+1}$. That 1 elements has to select from $A$. It’s easy to see $S_{a,b+1} = S_{a,b}\cup \{m\}$ or $S_{a,b+1} = S_{a,b}\cup \{n\}$, so there are 2 choice for $S_{a,b+1}$. Same thing applies to $S_{a+1,b}$.

 \textbf{Filling in the rest of the grid} We used our proved lemma, and we can fill in $S_{0,1}, S_{1,2} ... S_{n-1,n}$  then we can fill in the next diagonal, until all $S_{i,j}$ are filled, where $i \leq j$. But, we haven’t finished everything! Fortunately, filling out the rest of the diagonals in a similar fashion is pretty simple. And, it’s easy to see that we have made $n(n+1)$ decisions, each with 2 choices, when filling out the rest of the grid, so there are $2^{n(n+1)}$ ways to finish off.

 \textbf{Finishing off} To finish off, we have $\frac{(2n)!}{2^n} \cdot 2^{n(n+1)}$ ways to fill in the grid, which gets us $\boxed{(2n)! \cdot 2^{n^2}}$

 - AlexLikeMath

 


 


	\item (\textit{USAJMO 2019}) \textbf{Note:} For any geometry problem whose statement begins with an asterisk $(*)$, the first page of the solution must be a large, in-scale, clearly labeled diagram. Failure to meet this requirement will result in an automatic 1-point deduction.


 Two rational numbers $\frac{m}{n}$ and $\frac{n}{m}$ are written on a blackboard, where $m$ and $n$ are relatively prime positive integers. At any point, Evan may pick two of the numbers $x$ and $y$ written on the board and write either their arithmetic mean $\frac{x+y}{2}$ or their harmonic mean $\frac{2xy}{x+y}$ on the board as well. Find all pairs $(m,n)$ such that Evan can write $1$ on the board in finitely many steps.


 



\textbf{Solution}

We claim that all odd $m, n$ work if $m+n$ is a positive power of 2.

 Proof:We first prove that $m+n=2^k$ works. By weighted averages we have that $\frac{n(\frac{m}{n})+(2^k-n)\frac{n}{m}}{2^k}=\frac{m+n}{2^k}=1$ can be written, so the solution set does indeed work. We will now prove these are the only solutions.

 Assume that $m+n\ne 2^k$, so then $m+n\equiv 0\pmod{p}$ for some odd prime $p$. Then $m\equiv -n\pmod{p}$, so $\frac{m}{n}\equiv \frac{n}{m}\equiv -1\pmod{p}$. We see that the arithmetic mean is $\frac{-1+(-1)}{2}\equiv -1\pmod{p}$ and the harmonic mean is $\frac{2(-1)(-1)}{-1+(-1)}\equiv -1\pmod{p}$, so if 1 can be written then $1\equiv -1\pmod{p}$ and $2\equiv 0\pmod{p}$ which is obviously impossible, and we are done.

 -Stormersyle

 


\end{enumerate}